% TEMPLATE-MILC-v3.tex - plantilla maestra para todas las guias MILC v3
% Ver scripts/generadores/build-guia-milc.py para la lista completa de
% claves esperadas. El script valida que no queden placeholders sin
% reemplazar antes de escribir el archivo final.

\documentclass[11pt,letterpaper]{article}
\usepackage[margin=1.75cm]{geometry}
\usepackage[table]{xcolor}
\usepackage{fontspec}
\usepackage[spanish]{babel}
\usepackage{tikz}
\usepackage[most]{tcolorbox}
\usepackage{tabularx}
\usepackage{array}
\usepackage{fancyhdr}
\usepackage{titlesec}
\usepackage{ragged2e}
\usepackage{enumitem}
\usepackage{microtype}

\setmainfont{Helvetica Neue}
\setsansfont{Helvetica Neue}
\setlength{\parindent}{0pt}
\setlength{\parskip}{6pt}
\hyphenpenalty=200
\exhyphenpenalty=200
\pretolerance=400
\tolerance=2000
\setlength{\emergencystretch}{1em}
\renewcommand{\arraystretch}{1.25}
\setlist{leftmargin=5mm,itemsep=1mm,topsep=1mm}
\newcolumntype{Y}{>{\RaggedRight\arraybackslash}X}

\definecolor{milcMagenta}{HTML}{E6007E}
\definecolor{milcVino}{HTML}{5A0038}
\definecolor{milcTurquesa}{HTML}{0093A5}
\definecolor{milcVerde}{HTML}{58A923}
\definecolor{milcMostaza}{HTML}{E5B400}
\definecolor{milcOcre}{HTML}{B86600}
\definecolor{milcNegro}{HTML}{241816}
\definecolor{milcGris}{HTML}{F7F7F4}
\definecolor{milcLinea}{HTML}{E8E2DD}
\definecolor{milcGradePrimary}{HTML}{FF2D87}
\definecolor{milcGradeDark}{HTML}{9F1B56}
\definecolor{milcGradeSoft}{HTML}{FFE0EE}
\definecolor{milcGradeAccent}{HTML}{CFE7FF}
\definecolor{milcGradeText}{HTML}{FFFFFF}
\color{milcNegro}

\pagestyle{fancy}
\fancyhf{}
\lhead{\small Tecnología e Informática · Grado 10}
\rhead{\small Guía 27 · MILC v3}
\cfoot{\small\thepage}
\renewcommand{\headrulewidth}{0pt}

\newtcolorbox{titlebox}[1]{
  enhanced,colback=#1,colframe=#1,arc=7mm,boxrule=0pt,
  left=7mm,right=7mm,top=3mm,bottom=3mm,
  before skip=8pt,after skip=8pt,
  fontupper=\sffamily\bfseries\Large\color{white}
}

\newtcolorbox{softbox}[3]{
  enhanced,colback=#2,colframe=#1,borderline west={4pt}{0pt}{#1},
  arc=2mm,boxrule=.4pt,left=4mm,right=4mm,top=3mm,bottom=3mm,
  before skip=6pt,after skip=8pt,title={#3},coltitle=white,
  colbacktitle=#1,fonttitle=\sffamily\bfseries,fontupper=\RaggedRight,
  attach boxed title to top left={xshift=3mm,yshift=-2mm},
  boxed title style={arc=2mm, boxrule=0pt}
}

\newtcolorbox{drawbox}[2]{
  enhanced,colback=white,colframe=#1,arc=4mm,boxrule=1pt,
  left=5mm,right=5mm,top=4mm,bottom=4mm,before skip=8pt,after skip=8pt,
  title={#2},coltitle=white,colbacktitle=#1,fonttitle=\sffamily\bfseries,
  fontupper=\RaggedRight,
  attach boxed title to top left={xshift=4mm,yshift=-2mm},
  boxed title style={arc=3mm, boxrule=0pt}
}

\newtcolorbox{infoband}[1]{
  enhanced,colback=#1!8,colframe=#1!50,arc=2mm,boxrule=.5pt,
  left=4mm,right=4mm,top=2mm,bottom=2mm,before skip=4pt,after skip=4pt,
  fontupper=\small\RaggedRight
}

% Puente narrativo entre secciones (contrato editorial Conéctate).
% Da continuidad: "Acabas de X. Ahora vas a Y porque Z."
% Visualmente: banda izquierda en color del grado, texto en itálica pequeña.
\newtcolorbox{puentebox}{
  enhanced,colback=milcGradeSoft,colframe=milcGradeDark,
  borderline west={3pt}{0pt}{milcGradeDark},
  arc=0mm,boxrule=0pt,left=5mm,right=4mm,top=2.5mm,bottom=2.5mm,
  before skip=4pt,after skip=8pt,
  fontupper=\itshape\small\color{milcNegro}\RaggedRight
}
\newcommand{\puente}[1]{%
  \begin{puentebox}%
    \textcolor{milcGradeDark}{\textbf{\textrightarrow}}\hspace{2mm}#1%
  \end{puentebox}%
}

\newcommand{\linead}[1]{\noindent\color{milcLinea}\rule{#1}{0.5pt}\color{milcNegro}}
\newcommand{\respuesta}[1]{\par\vspace{2mm}\foreach \n in {1,...,#1}{\noindent\color{milcLinea}\rule{\linewidth}{0.5pt}\par\vspace{5mm}}\color{milcNegro}}
\newcommand{\checkbox}{\textcolor{milcGradeDark}{\fbox{\phantom{x}}}\hspace{5pt}}

\newcommand{\phasecard}[4]{
\begin{tcolorbox}[enhanced,colback=#3,colframe=#2,arc=3mm,boxrule=.5pt,height=3.05cm,valign=center,halign=center,left=1.5mm,right=1.5mm,top=2mm,bottom=2mm]
{\sffamily\bfseries\fontsize{22}{24}\selectfont\textcolor{#2}{#1}}\par
{\sffamily\bfseries\small #4}
\end{tcolorbox}
}

\begin{document}
\thispagestyle{empty}
\begin{tikzpicture}[remember picture,overlay]
  \fill[white] (current page.south west) rectangle (current page.north east);
  \path[fill=milcGradePrimary,rounded corners=16mm]
    ([xshift=.55cm,yshift=-.55cm]current page.north west) rectangle
    ([xshift=-.55cm,yshift=3.65cm]current page.south east);

  \node[anchor=north west,text=milcGradeText] at ([xshift=2.0cm,yshift=-1.85cm]current page.north west)
    {\sffamily\bfseries\fontsize{14}{16}\selectfont GRADO};
  \node[anchor=north west,text=milcGradeText,opacity=.82] at ([xshift=2.0cm,yshift=-2.75cm]current page.north west)
    {\sffamily\bfseries\fontsize{9}{11}\selectfont SERIE GUÍAS · TECNOLOGÍA E INFORMÁTICA};

  \begin{scope}[shift={([xshift=-3.05cm,yshift=-2.55cm]current page.north east)}]
    \fill[white,opacity=.80] (-.62,-.52) rectangle (.62,.52);
    \draw[milcGradeText,line width=1pt] (-.62,-.52) rectangle (.62,.52);
    \fill[milcGradeDark] (-.42,-.38) rectangle (-.22,.06);
    \fill[milcGradeAccent] (-.10,-.38) rectangle (.10,.24);
    \fill[milcGradePrimary!60!black] (.22,-.38) rectangle (.42,.38);
  \end{scope}

  \node[anchor=west,text=milcGradeText] at ([xshift=1.9cm,yshift=-12.85cm]current page.north west)
    {\sffamily\bfseries\fontsize{124}{124}\selectfont 10};
  \draw[milcGradeText,line width=1.7pt]
    ([xshift=4.52cm,yshift=-11.68cm]current page.north west) circle (.13cm);

  \node[anchor=west,text=milcGradeText,text width=8.7cm,align=left] at ([xshift=10.35cm,yshift=-11.6cm]current page.north west)
    {
     {\sffamily\bfseries\fontsize{19}{22}\selectfont Plan de negocio ---\\Business Model Canvas\\asistido por IA}\\[4mm]
     {\sffamily\bfseries\fontsize{23}{26}\selectfont Guía 27-10-TIC}\\[2mm]
     {\sffamily\fontsize{13}{16}\selectfont Período 3 · Ofimática con IA: contabilidad, datos y emprendimiento}\\[5mm]
     {\sffamily\bfseries\fontsize{11}{13}\selectfont Institución Educativa Sor María Juliana}\\[1mm]
     {\sffamily\fontsize{10}{12}\selectfont Autor: Dr. Álvaro Cárdenas Orozco}
    };

  \path[fill=milcGradeSoft,rounded corners=7mm]
    ([xshift=1.35cm,yshift=.85cm]current page.south west) rectangle
    ([xshift=-1.35cm,yshift=4.25cm]current page.south east);
  \draw[milcGradeDark,line width=.65pt,rounded corners=7mm]
    ([xshift=1.35cm,yshift=.85cm]current page.south west) rectangle
    ([xshift=-1.35cm,yshift=4.25cm]current page.south east);
  \node[anchor=center,text=milcNegro,text width=17.0cm,align=center] at ([yshift=2.55cm]current page.south)
    {\sffamily\fontsize{10}{12}\selectfont Metodología MILC · Apertura ancestral · Pensamiento computacional · Triángulo Dussel-Estoicismo-Floridi\\
    \textcolor{milcGradeDark}{\bfseries Producto:} Business Model Canvas completo del microemprendimiento con los 9 bloques llenos (segmento de clientes, propuesta de valor, canales, relaciones con clientes, fuentes de ingresos, recursos clave, actividades clave, socios clave, estructura de costos) + asistencia de IA para refinar al menos 3 bloques específicos + 3 supuestos críticos identificados que requerirán verificación con MVP en futuro + dibujo o documento del Canvas en formato visual (Canva, hoja apaisada, plantilla online gratuita)};
\end{tikzpicture}
\mbox{}
\newpage

\section*{Apertura: Plan de negocio con IA --- Business Model Canvas asistido}

\begin{softbox}{milcGradeDark}{milcGradeSoft}{Saber ancestral}
En la galería del mercado del centro de Cartago, en las plazas comerciales de Cali y Buenaventura, en los pueblos del Pacífico colombiano, hubo durante siglos un oficio que cualquier comerciante del barrio practicaba sin llamarlo \emph{modelo de negocio}: \textbf{el trueque}. La señora traía racimo de plátano y se iba con pescado; el pescador llegaba con corvinas y volvía con queso; el quesero bajaba con cuajada y subía con frutas. Ese intercambio ancestral tenía una arquitectura silenciosa que cualquier emprendedor moderno reconocería: había un \textbf{segmento} (quién intercambiaba con quién), una \textbf{propuesta de valor} (qué daba cada uno), un \textbf{canal} (la plaza, el día, la hora), una \textbf{relación} (la confianza acumulada de años), una \textbf{moneda implícita} (volumen, calidad, frescura) y unos \textbf{recursos clave} (terreno, semilla, embarcación, conocimiento). El \textbf{Business Model Canvas} moderno no inventa nada nuevo: \textbf{ordena en 9 casillas lo que el trueque hacía intuitivamente}. La diferencia es la escala y la urgencia de hacerlo explícito. La IA hoy ayuda a refinar el Canvas, pero el criterio sigue siendo del emprendedor.
\end{softbox}

\begin{softbox}{milcTurquesa}{milcGris}{Saber contemporáneo}
El \textbf{Business Model Canvas} (BMC) fue creado por Alex Osterwalder en 2010 y es la herramienta visual más usada para diseñar modelos de negocio. Sus \textbf{9 bloques} se organizan en una hoja apaisada: (1) \textbf{Segmento de clientes}: para quién creas valor. (2) \textbf{Propuesta de valor}: qué problema resuelves o qué deseo satisfaces. (3) \textbf{Canales}: cómo entregas el valor (físico, digital, mixto). (4) \textbf{Relaciones con clientes}: cómo te relacionas con ellos (personal, automatizado, comunidad). (5) \textbf{Fuentes de ingresos}: de dónde sale el dinero (venta única, suscripción, comisión, etc.). (6) \textbf{Recursos clave}: qué necesitas tener (físico, intelectual, humano, financiero). (7) \textbf{Actividades clave}: qué haces para entregar el valor. (8) \textbf{Socios clave}: quién colabora contigo. (9) \textbf{Estructura de costos}: qué te cuesta operar. La regla profesional: \textbf{el Canvas se llena con datos verificables, no con suposiciones}. Los hallazgos del estudio de mercado (S6) alimentan al menos 3 bloques: segmento, propuesta de valor, fuentes de ingresos. La \textbf{IA como copiloto} ayuda a refinar bloques específicos con prompts como: \emph{``Actúa como mentor de emprendimiento. Aquí está el bloque X de mi Canvas: [pegar]. Sugiere 3 mejoras concretas''}.
\end{softbox}

\begin{softbox}{milcMostaza}{milcGris}{Pregunta-puente}
\textit{¿Qué sabía el comerciante del trueque al tener clara la arquitectura silenciosa del intercambio, que el emprendedor novato olvida cuando arma negocio sin estructura? ¿Y por qué los \emph{supuestos críticos} del Canvas son más importantes que los bloques mismos?}
\end{softbox}

\begin{softbox}{milcVerde}{milcGris}{Saber y saber hacer}
Al terminar podrás: (1) \textbf{identificar} los 9 bloques del BMC y qué responde cada uno; (2) \textbf{analizar} qué bloques se sostienen en datos del estudio de mercado y cuáles aún son suposiciones; (3) \textbf{crear} el Canvas completo con asistencia de IA en al menos 3 bloques; (4) \textbf{evaluar} qué 3 supuestos críticos requerirán verificación con MVP futuro.
\end{softbox}

\small
\begin{tabularx}{\textwidth}{>{\bfseries}p{3.0cm}Y}
\rowcolor{milcGradePrimary!12}Autor & Dr. Álvaro Cárdenas Orozco\\
DBA & Aplicación de ofimática asistida por IA a contabilidad básica y emprendimiento (MEN, T\&I)\\
Referentes & Lean Startup · Phronesis financiera · Ética del dato empresarial\\
Producto final & Business Model Canvas completo del microemprendimiento con los 9 bloques llenos (segmento de clientes, propuesta de valor, canales, relaciones con clientes, fuentes de ingresos, recursos clave, actividades clave, socios clave, estructura de costos) + asistencia de IA para refinar al menos 3 bloques específicos + 3 supuestos críticos identificados que requerirán verificación con MVP en futuro + dibujo o documento del Canvas en formato visual (Canva, hoja apaisada, plantilla online gratuita)\\
\end{tabularx}
\normalsize

\puente{Antes de empezar, mira el viaje completo. En la sesión 6 hiciste estudio de mercado real. Hoy lo formalizas en el \textbf{lienzo del negocio}: 9 bloques que describen el modelo completo. Las próximas sesiones (comunicación empresarial, proyecto integrador, sustentación) usarán el Canvas como base de la propuesta final.}

\begin{titlebox}{milcGradeDark}
Ruta MILC: así vamos a aprender
\end{titlebox}
\begin{tabularx}{\textwidth}{>{\centering\arraybackslash}X>{\centering\arraybackslash}X>{\centering\arraybackslash}X>{\centering\arraybackslash}X}
\phasecard{1}{milcVerde}{milcGris}{Escuta\\[-1mm]\small Observo, escucho y pregunto.} &
\phasecard{2}{milcTurquesa}{milcGris}{Sistematización\\[-1mm]\small Pensamiento computacional.} &
\phasecard{3}{milcMagenta}{milcGris}{Praxis\\[-1mm]\small Construyo y aplico.} &
\phasecard{4}{milcVino}{milcGris}{Evaluación\\[-1mm]\small Reflexión liberadora.}
\end{tabularx}


\puente{Antes de teorizar, vas a hacer un \emph{ejercicio del trueque}: pensar en el microemprendimiento de S4-S6 y responder en 1 frase cada una de las 9 preguntas del Canvas. Sin profundizar; solo identificar.}

\begin{titlebox}{milcVerde}
Fase 1 · Escuta
\end{titlebox}

\begin{softbox}{milcVerde}{milcGris}{Escena para escuchar}
\textbf{Actividad 1 · IDENTIFICA --- Versión cruda del Canvas en 9 frases} (15 min · individual). Toma el microemprendimiento de S4-S6. Responde en 1 frase cada bloque del Canvas: (1) Segmento: ¿quién es mi cliente? (2) Propuesta de valor: ¿qué problema le resuelvo? (3) Canales: ¿cómo le entrego? (4) Relaciones: ¿cómo me relaciono con él? (5) Ingresos: ¿cómo me paga? (6) Recursos: ¿qué necesito tener? (7) Actividades: ¿qué hago? (8) Socios: ¿quién me ayuda? (9) Costos: ¿qué me cuesta? Sin profundizar; solo identificar.
\end{softbox}

\checkbox Respondí las 9 preguntas en 1 frase cada una.\\
\checkbox Identifiqué dónde tengo claridad y dónde tengo dudas.\\
\checkbox Reconozco que algunos bloques se sostienen en datos y otros aún no.

\begin{infoband}{milcVerde}
\textbf{Lo que va al cuaderno (Actividad 1 · IDENTIFICA).} Título: «Actividad 1 --- Canvas en 9 frases». \textbf{Formato:} 9 frases identificadas + marca de cuáles se sostienen en datos de S6 y cuáles aún son suposiciones. \textbf{Sabes que terminaste cuando} ves qué bloques necesitan más trabajo.
\end{infoband}

\puente{Ya tienes versión cruda. Ahora vas a entender la \textbf{anatomía profesional del Business Model Canvas}: 9 bloques, conexión con estudio de mercado, los 5 errores típicos.}

\begin{titlebox}{milcTurquesa}
Fase 2 · Sistematización con pensamiento computacional
\end{titlebox}

\begin{softbox}{milcTurquesa}{milcGris}{Cuatro pilares aplicados al tema}
Tu Canvas crudo te muestra el panorama. Ahora necesitas la \textbf{anatomía profesional} de los 9 bloques y la conexión con datos reales.
\end{softbox}

\small
\begin{tabularx}{\textwidth}{>{\bfseries\RaggedRight\arraybackslash}p{3.6cm}Y}
\rowcolor{milcTurquesa!90}\color{white}Pilar & \color{white}\bfseries Aplicación al tema\\
1. Descomposición & Los \textbf{9 bloques del Canvas} se descomponen así: (1) \textbf{Segmento de clientes}: \emph{¿para quién creo valor?}. Pueden ser 1-3 segmentos distintos. (2) \textbf{Propuesta de valor}: \emph{¿qué problema resuelvo, qué beneficio entrego?}. La pieza más importante del Canvas. (3) \textbf{Canales}: \emph{¿cómo llego al cliente?} (físico, digital, mixto). (4) \textbf{Relaciones con clientes}: \emph{¿cómo me relaciono?} (personal, automatizado, comunidad, autoservicio). (5) \textbf{Fuentes de ingresos}: \emph{¿cómo cobro?} (venta única, suscripción, comisión, freemium). (6) \textbf{Recursos clave}: \emph{¿qué necesito?} (físico: local; intelectual: receta, marca; humano: personal; financiero: capital). (7) \textbf{Actividades clave}: \emph{¿qué hago?} (producción, marketing, atención). (8) \textbf{Socios clave}: \emph{¿quién me ayuda?} (proveedores, aliados, distribuidores). (9) \textbf{Estructura de costos}: \emph{¿qué me cuesta?} (fijos + variables, sesión 4).\\
2. Reconocimiento de patrones & La \textbf{conexión con estudio de mercado} (S6) es crucial. Los hallazgos numéricos alimentan al menos 3 bloques: (a) Segmento se afina con la demografía de los encuestados que mostraron interés. (b) Propuesta de valor se ajusta con lo que dijeron que necesitan. (c) Fuentes de ingresos se valida con el precio que dijeron pagar. El resto de bloques se construye con criterio del emprendedor + asistencia de IA.\\
3. Abstracción & Los \textbf{supuestos críticos} son los bloques cuyo error rompería todo el modelo. Típicamente: (a) \textbf{Propuesta de valor}: supuesto de que el cliente quiere lo que ofrezco. (b) \textbf{Canales}: supuesto de que sé dónde encontrarlo. (c) \textbf{Fuentes de ingresos}: supuesto de que pagará lo que pido. Estos 3 son los blancos preferidos del MVP que verificarías en S9.\\
4. Algoritmo conceptual & Algoritmo: (1) tomar Canvas crudo de Actividad 1. (2) refinar bloques con datos de S6. (3) usar IA para refinar 3 bloques específicos. (4) crear versión visual en Canva, hoja apaisada o plantilla online gratuita (canvanizer.com es gratis). (5) identificar 3 supuestos críticos a verificar con MVP. (6) verificar coherencia entre bloques.\\
\end{tabularx}
\normalsize

\begin{drawbox}{milcTurquesa}{Anatomía del BMC}
\textbf{9 bloques:} segmento / propuesta / canales / relaciones / ingresos / recursos / actividades / socios / costos. \\[1mm] \textbf{Datos:} S6 alimenta 3 bloques principales. \\[1mm] \textbf{IA:} refina 3 bloques específicos. \\[1mm] \textbf{Supuestos críticos:} 3 que requieren verificación.
\end{drawbox}

\begin{softbox}{milcMostaza}{milcGris}{Errores comunes y su corrección}
(1) \textbf{Segmento muy amplio}: \emph{``todo el mundo''} es ningún cliente. (2) \textbf{Propuesta de valor genérica}: \emph{``mejor calidad''} no es propuesta. (3) \textbf{Ingresos sin sustento}: precio basado en \emph{``creo que es justo''} sin estudio. (4) \textbf{Sin identificar supuestos críticos}: el Canvas se ve completo pero no se sabe qué validar. (5) \textbf{Aceptar todo lo que la IA propone}: la IA puede sugerir bloques irrealistas para tu contexto.
\end{softbox}

\begin{infoband}{milcTurquesa}
\textbf{Lo que va al cuaderno (Actividad 2 · ANALIZA).} Título: «Actividad 2 --- Anatomía BMC». \textbf{Formato:} (a) los 9 bloques con su pregunta; (b) conexión con S6; (c) los 5 errores con marca. \textbf{Sabes que terminaste cuando} podrías auditar Canvas de otro proyecto.
\end{infoband}

\puente{Tienes el método. Toca aplicarlo: armas el Canvas completo en formato visual, usas IA para refinar 3 bloques específicos, identificas 3 supuestos críticos.}

\begin{titlebox}{milcMagenta}
Fase 3 · Praxis con pensamiento computacional
\end{titlebox}

\begin{softbox}{milcMagenta}{milcGris}{Construyendo el producto con los cuatro pilares}
\textbf{Actividad 3 · CREA + EVALÚA --- Canvas completo + IA + supuestos críticos} (40 min · individual con dispositivo). Hora de aplicar. Construyes Canvas visual con los 9 bloques, usas IA para refinar 3 específicos, identificas 3 supuestos críticos.
\end{softbox}

\small
\begin{tabularx}{\textwidth}{>{\bfseries\RaggedRight\arraybackslash}p{3.6cm}Y}
\rowcolor{milcMagenta!90}\color{white}Pilar & \color{white}\bfseries Aplicación al producto\\
Descomposición & Tu Canvas se construye en 6 pasos: (1) abrir Canva con plantilla \emph{Business Model Canvas} (gratis), hoja apaisada en Docs/LibreOffice, o canvanizer.com (gratuito en línea). (2) llenar los 9 bloques con frases concisas (3-7 palabras por idea), conectando con los hallazgos del estudio de mercado. (3) elegir 3 bloques específicos para refinar con IA. (4) escribir prompt para cada uno: \emph{``Actúa como mentor de emprendimiento. Mi microemprendimiento es [X]. El bloque \emph{[nombre]} de mi Canvas actualmente dice [pegar]. Sugiere 3 mejoras concretas basadas en lo que dirían 17 de 20 clientes que entrevisté que [pegar hallazgos]''}. (5) integrar las mejores sugerencias de la IA. (6) identificar 3 supuestos críticos que requerirán MVP en futuro.\\
Patrones & Sigue el patrón \textbf{``datos sobre suposiciones''}: cada bloque debe poder responder a la pregunta \emph{``¿en qué dato me baso para afirmar esto?''}. Si la respuesta es \emph{``en mi intuición''}, hay supuesto crítico que verificar. Si la respuesta es \emph{``en el hallazgo de mi estudio''}, el bloque tiene sustento.\\
Abstracción & Los \textbf{3 supuestos críticos} se identifican así: (a) ¿qué afirmación del Canvas, si fuera falsa, haría caer el modelo entero? (b) ¿qué supuesto no está verificado todavía con datos? (c) ¿qué supuesto sería el más caro de equivocarse? Esos 3 supuestos serán los blancos del MVP de la sesión 9 (proyecto integrador).\\
Algoritmo de armado & Algoritmo: (1) Canvas en formato visual con 9 bloques. (2) refinar con IA 3 bloques. (3) identificar 3 supuestos críticos. (4) verificar coherencia entre bloques. (5) versión final del Canvas guardada para compartir.\\
\end{tabularx}
\normalsize

\begin{drawbox}{milcMagenta}{Checklist antes de entregar}
\checkbox 9 bloques llenos en formato visual. \\[1mm] \checkbox 3 bloques refinados con IA documentada. \\[1mm] \checkbox Datos del estudio de mercado conectados al Canvas. \\[1mm] \checkbox 3 supuestos críticos identificados. \\[1mm] \checkbox Coherencia entre bloques verificada. \\[1mm] \checkbox Canvas guardado y compartible.
\end{drawbox}

\begin{softbox}{milcMagenta}{milcGris}{Plantilla de trabajo}
\textbf{9 bloques.} \textit{Cada uno con frase concisa.} \\[1mm] \textbf{3 bloques con IA.} \textit{Refinados con prompts específicos.} \\[1mm] \textbf{Conexión datos.} \textit{Hallazgos S6 alimentando Canvas.} \\[1mm] \textbf{3 supuestos críticos.} \textit{A verificar en MVP futuro.}
\end{softbox}

\begin{infoband}{milcMagenta}
\textbf{Lo que va al cuaderno (Actividad 3 · CREA + EVALÚA).} Título: «Actividad 3 --- Mi Canvas con IA». \textbf{Formato:} (a) imagen o link del Canvas; (b) prompts usados con IA; (c) 3 supuestos críticos identificados. \textbf{Sabes que terminaste cuando} alguien viendo tu Canvas entiende el modelo en 5 minutos.
\end{infoband}

\puente{Lo que sigue resume \emph{qué} entregas hoy: Canvas visual + IA aplicada a 3 bloques + 3 supuestos críticos. El criterio principal: que el Canvas refleje datos reales del estudio de mercado, no suposiciones del emprendedor.}

\begin{titlebox}{milcMostaza}
Producto de la sesión
\end{titlebox}
\textbf{BMC completo + 3 bloques refinados con IA + 3 supuestos críticos}.
\begin{softbox}{milcMostaza}{milcGris}{Criterios de calidad}
(1) 9 bloques llenos con frases concisas. (2) Datos del estudio de mercado integrados. (3) 3 bloques refinados con IA documentada. (4) 3 supuestos críticos identificados. (5) Coherencia entre bloques. (6) Canvas visual compartible.
\end{softbox}

\puente{Antes de cerrar, mira el Canvas desde las cinco dimensiones humanas. Diseñar el modelo con claridad es respeto por uno mismo y por futuros socios; improvisar el negocio es invitar al fracaso con buenos modales.}

\begin{titlebox}{milcVino}
Fase 4 · Evaluación liberadora — cinco dimensiones
\end{titlebox}

\begin{softbox}{milcVino}{milcGris}{Sentido de la evaluación}
Evaluar no es solo revisar una nota. Es reconocer qué aprendí, qué cambió en mí, qué emociones manejé, cómo me relaciono con la ciudad y la comunidad, y qué le dejaré a quien viene detrás.
\end{softbox}

\textbf{Mi reflexión liberadora en cinco dimensiones.} Completa cada frase iniciada:

\small
\begin{tabularx}{\textwidth}{>{\bfseries\RaggedRight\arraybackslash}p{3.4cm}Y>{\RaggedRight\arraybackslash}p{4.6cm}}
\rowcolor{milcVino}\color{white}Dimensión & \color{white}\bfseries Mi respuesta (frame starter) & \color{white}\bfseries Algo que descubrí\\
1. Personal & Lo que más cambió en mí fue \linead{6.5cm} & \linead{4.2cm}\\
2. Emocional & La emoción más fuerte fue \linead{2.5cm} y la regulé \linead{3cm} & \linead{4.2cm}\\
3. Ciudadana & En democracia esto importa porque \linead{6cm} & \linead{4.2cm}\\
4. Local (Cartago) & En mi barrio sería útil para \linead{6.5cm} & \linead{4.2cm}\\
5. Intergeneracional & A mi abuela/abuelo le diría \linead{6.5cm} & \linead{4.2cm}\\
\end{tabularx}
\normalsize

\begin{infoband}{milcVino}
\textbf{Para la próxima sustentación.} Vuelve a esta tabla en seis meses. Verás que las respuestas cambian — no porque lo aprendido cambie, sino porque tú cambias. La evaluación liberadora es una conversación contigo a través del tiempo.
\end{infoband}


\puente{Y para terminar, tres voces sobre la arquitectura del negocio. Dussel sobre los modelos que sirven al cliente, Marco Aurelio sobre la disciplina de los 9 bloques, Floridi sobre el Canvas como instrumento de claridad estratégica.}

\begin{titlebox}{milcGradeDark}
Triángulo de pensamiento · cierre filosófico
\end{titlebox}

\begin{softbox}{milcVino}{milcGris}{Dussel · lente del nosotros}
\textit{``Un modelo de negocio que sirve al cliente sin extraer es liberador; uno que solo busca lucro reproduce la lógica colonial.''}

Tu Canvas puede diseñarse de 2 maneras: (a) servir al cliente entregando valor proporcional a lo que cobra. (b) extraer del cliente más de lo que entrega. La filosofía de la liberación pide la primera. Las microempresas del barrio que duran 30 años son las que servían honestamente; las que se quebraban rápido eran las que extraían sin entregar. La phronesis emprendedora consiste en \textbf{diseñar modelos que sirven, no que solo extraen}.

\textbf{¿Mi Canvas diseña negocio que sirve al cliente, o solo busca extraer valor de él?}
\end{softbox}
\linead{\linewidth}\\[1mm]
\linead{\linewidth}

\begin{softbox}{milcMostaza}{milcGris}{Estoicismo · Marco Aurelio · cuidado interior}
\textit{``Llenar los 9 bloques con disciplina es virtud emprendedora; saltarse alguno es debilidad que se descubre al lanzar.''}

Es tentador saltarse bloques que \emph{``ya pensaré después''} (socios, costos). La sabiduría estoica recomienda lo opuesto: \emph{llena los 9, aunque algunos sean preliminares}. La incompletitud del Canvas se traduce en sorpresas costosas al lanzar.

\textbf{¿Mis 9 bloques están todos llenos, o salté alguno por comodidad?}
\end{softbox}
\linead{\linewidth}\\[1mm]
\linead{\linewidth}

\begin{softbox}{milcTurquesa}{milcGris}{Floridi · lente de la infoesfera}
\textit{``El Canvas con IA como copiloto es la nueva práctica del emprendimiento del siglo XXI.''}

Tradicionalmente, el Canvas se llenaba a mano o con mentor humano. Hoy la IA puede asistir refinando bloques específicos. La phronesis emprendedora contemporánea consiste en \textbf{usar la IA donde acelera, conservar el criterio donde decide}. La decisión final sobre el modelo es del emprendedor, no de la IA.

\textbf{¿Estoy usando IA como copiloto del Canvas, o le delego la decisión sobre el modelo?}
\end{softbox}
\linead{\linewidth}\\[1mm]
\linead{\linewidth}

\begin{titlebox}{milcVino}
Compromiso final
\end{titlebox}

\small
\begin{tabularx}{\textwidth}{>{\RaggedRight\arraybackslash}p{3.0cm}YYY}
\rowcolor{milcVino}\color{white}\bfseries Criterio & \color{white}\bfseries Lo logré & \color{white}\bfseries En proceso & \color{white}\bfseries Necesito apoyo\\
Comprensión del tema & Explico ideas centrales con mis palabras. & Reconozco ideas pero falta precisión. & Necesito repasar conceptos base.\\
Pensamiento computacional & Apliqué los 4 pilares al diseñar mi producto. & Apliqué algunos pilares. & Necesito releer la fase 2.\\
Aplicación y producto & Mi producto cumple los criterios y se puede revisar. & Producto funcional con ajustes pendientes. & Aún en construcción.\\
Reflexión 5 dimensiones & Respondí cada dimensión con honestidad. & Respondí algunas con detalle. & Mis respuestas son muy generales.\\
Triángulo de pensamiento & Conecté las 3 voces con mi trabajo real. & Conecté al menos una. & No relaciono las voces con mi proceso.\\
\end{tabularx}
\normalsize

\textbf{Hoy entendí que:}\\[1mm]
\linead{\linewidth}\\[1mm]
\linead{\linewidth}

\textbf{Mi compromiso para la próxima sesión es:}\\[1mm]
\linead{\linewidth}\\[1mm]
\linead{\linewidth}

\begin{softbox}{milcMostaza}{milcGris}{Cita de despedida · Marco Aurelio}
\textit{``El obstáculo en el camino se vuelve el camino. Lo que estorba al camino se vuelve camino.''}\\[1mm]
Cada error de hoy, bien observado, se convierte en pieza del camino que pisarás mañana.
\end{softbox}

\small
\begin{tabularx}{\textwidth}{>{\bfseries}p{3.6cm}Y}
\rowcolor{milcGradeDark!12}Nombre completo: & \linead{10cm}\\
Fecha: & \linead{4cm} \quad Curso/grupo: \linead{4cm}\\
Mi firma: & \linead{9cm}\\
\end{tabularx}
\normalsize

\begin{infoband}{milcGradeDark}
\textbf{Para la próxima sesión.} Llega con tu Canvas visual y los 3 supuestos críticos identificados. En la sesión 8 vas a aprender \textbf{comunicación empresarial profesional con IA}: cómo presentar el Canvas y la propuesta a clientes y aliados. El Canvas de hoy es la arquitectura; la comunicación de mañana es cómo se cuenta.
\end{infoband}

\end{document}
